\documentclass[11pt,a4paper]{article}
\usepackage[utf8]{inputenc}
\usepackage[T1]{fontenc}
\usepackage[portuguese]{babel}
\usepackage{geometry}
\geometry{margin=2.5cm}
\usepackage{listings}
\usepackage{xcolor}
\usepackage{amsmath}
\usepackage{amssymb}
\usepackage{hyperref}
\usepackage{enumitem}
\usepackage{tcolorbox}
\usepackage{tabularx}
\usepackage{booktabs}

\definecolor{codebg}{rgb}{0.95,0.95,0.95}
\definecolor{codegreen}{rgb}{0.1,0.5,0.1}
\definecolor{codepurple}{rgb}{0.5,0.1,0.5}
\definecolor{codegray}{rgb}{0.5,0.5,0.5}

\lstdefinestyle{python}{
    backgroundcolor=\color{codebg},
    commentstyle=\color{codegray}\itshape,
    keywordstyle=\color{codepurple}\bfseries,
    stringstyle=\color{codegreen},
    basicstyle=\ttfamily\small,
    breaklines=true,
    frame=single,
    language=Python,
    showstringspaces=false,
    tabsize=4,
    literate={á}{{\'a}}1 {ã}{{\~a}}1 {é}{{\'e}}1 {í}{{\'i}}1 {ó}{{\'o}}1 {õ}{{\~o}}1 {ú}{{\'u}}1 {ç}{{\c{c}}}1 {Á}{{\'A}}1 {Ã}{{\~A}}1 {É}{{\'E}}1 {Í}{{\'I}}1 {Ó}{{\'O}}1 {Õ}{{\~O}}1 {Ú}{{\'U}}1 {Ç}{{\c{C}}}1 {â}{{\^a}}1 {ê}{{\^e}}1 {ô}{{\^o}}1 {û}{{\^u}}1 {Â}{{\^A}}1 {Ê}{{\^E}}1 {Ô}{{\^O}}1 {Û}{{\^U}}1 {à}{{\`a}}1 {è}{{\`e}}1 {ì}{{\`i}}1 {ò}{{\`o}}1 {ù}{{\`u}}1
}
\lstset{style=python}

\newtcolorbox{solucao}{
  colback=yellow!5!white,
  colframe=orange!60!black,
  title=Solução proposta,
  fonttitle=\bfseries
}

\title{\textbf{Data Analysis with Python 2024--2025}\\Parte 5: Machine Learning (Regressão Linear)\\\large Soluções dos Exercícios}
\author{Soluções independentes baseadas nos enunciados originais}
\date{}

\begin{document}
\maketitle

\section*{Nota Importante}
Este material é uma adaptação das soluções independentes para os exercícios baseados no curso \textit{Data Analysis with Python 2024-2025}, oferecido pela \textbf{The University of Helsinki MOOC Center}. As soluções abaixo não são o gabarito oficial do curso.

\tableofcontents
\newpage

\section{Soluções - Regressão Linear}

A seguir estão as soluções propostas para os exercícios sobre Machine Learning e Regressão Linear.

\subsection*{Exercício 5.10 -- Regressão Linear}
\begin{solucao}
\begin{lstlisting}
import numpy as np
import matplotlib.pyplot as plt
from sklearn.linear_model import LinearRegression

def fit_line(x, y):
    model = LinearRegression(fit_intercept=True)
    # Adicionando um novo eixo para converter array 1D para 2D exigido pelo sklearn
    model.fit(x[:, np.newaxis], y)
    return model.coef_[0], model.intercept_

def main():
    x = np.array([1, 2, 3, 4, 5])
    y = np.array([2.5, 3.1, 4.8, 5.2, 6.7])
    
    slope, intercept = fit_line(x, y)
    print(f"Slope: {slope}")
    print(f"Intercept: {intercept}")
    
    plt.plot(x, y, 'o', label='Pontos originais')
    plt.plot(x, slope * x + intercept, color='red', label='Linha ajustada')
    plt.legend()
    plt.show()

if __name__ == "__main__":
    main()
\end{lstlisting}
Esta solução instancia um objeto \texttt{LinearRegression} para calcular a linha que melhor se ajusta aos dados. Em seguida, extraímos o primeiro valor dos coeficientes gerados e o intercepto para retorná-los como tupla.
\end{solucao}

\subsection*{Exercício 5.11 -- Dados Misteriosos (Mystery Data)}
\begin{solucao}
\begin{lstlisting}
import pandas as pd
from sklearn.linear_model import LinearRegression

def mystery_data():
    df = pd.read_csv("mystery_data.tsv", sep="\t")
    # As 5 primeiras colunas sao os atributos (features) e a ultima e a resposta (Y)
    X = df.iloc[:, 0:5]
    y = df.iloc[:, 5]
    
    model = LinearRegression(fit_intercept=False)
    model.fit(X, y)
    return model.coef_

def main():
    coefficients = mystery_data()
    for i, coef in enumerate(coefficients, 1):
        print(f"Coefficient of X{i} is {coef}")

if __name__ == "__main__":
    main()
\end{lstlisting}
O parâmetro \texttt{sep="\textbackslash t"} é utilizado na leitura do arquivo TSV. Dividimos o DataFrame nas variáveis explicativas (\texttt{X}) e na variável dependente (\texttt{y}). Foi configurado \texttt{fit\_intercept=False} de acordo com as instruções.
\end{solucao}

\subsection*{Exercício 5.12 -- Coeficiente de Determinação (R2)}
\begin{solucao}
\begin{lstlisting}
import pandas as pd
from sklearn.linear_model import LinearRegression

def coefficient_of_determination():
    df = pd.read_csv("mystery_data.tsv", sep="\t")
    X = df.iloc[:, 0:5]
    y = df.iloc[:, 5]
    
    model = LinearRegression(fit_intercept=True)
    
    # Avaliando todas as features juntas
    model.fit(X, y)
    scores = [model.score(X, y)]
    
    # Avaliando as features isoladamente
    for i in range(5):
        X_single = X.iloc[:, i].values.reshape(-1, 1)
        model.fit(X_single, y)
        scores.append(model.score(X_single, y))
        
    return scores

def main():
    scores = coefficient_of_determination()
    print(f"R2-score with feature(s) X: {scores[0]}")
    for i, score in enumerate(scores[1:], 1):
        print(f"R2-score with feature(s) X{i}: {score}")

if __name__ == "__main__":
    main()
\end{lstlisting}
Para treinar as features de modo avulso, é preciso isolá-las iterando sob as colunas do conjunto \texttt{X} e alterando o formato utilizando \texttt{reshape(-1, 1)}, um requisito para arrays únicos no Scikit-learn.
\end{solucao}

\subsection*{Exercício 5.13 -- O Clima de Ciclismo Continua}
\begin{solucao}
\begin{lstlisting}
import pandas as pd
from sklearn.linear_model import LinearRegression

def cycling_weather_continues(station):
    # Lendo os dados
    weather = pd.read_csv("src/weather_data.csv")
    cycling = pd.read_csv("src/cycling_data.csv")
    
    # Restringindo os dados ao ano de 2017
    cycling = cycling[cycling['Year'] == 2017]
    
    # Agrupando e somando os contadores diarios de ciclismo
    cycling_daily = cycling.groupby(['Year', 'Month', 'Day']).sum().reset_index()
    
    # Mesclando os dados por Ano, Mes e Dia
    merged = pd.merge(weather, cycling_daily, on=['Year', 'Month', 'Day'], how='inner')
    
    # Preenchendo os espacos vazios usando "forward fill"
    merged = merged.fillna(method='ffill')
    
    # Separando variaveis dependentes e independentes
    X = merged[['Precipitation amount (mm)', 'Snow depth (cm)', 'Air temperature (degC)']]
    y = merged[station]
    
    # Treinando o modelo
    model = LinearRegression(fit_intercept=True)
    model.fit(X, y)
    
    return model.coef_, model.score(X, y)

def main():
    station = 'Baana'
    coefs, score = cycling_weather_continues(station)
    
    print(f"Measuring station: {station}")
    print(f"Regression coefficient for variable 'precipitation': {coefs[0]:.1f}")
    print(f"Regression coefficient for variable 'snow depth': {coefs[1]:.1f}")
    print(f"Regression coefficient for variable 'temperature': {coefs[2]:.1f}")
    print(f"Score: {score:.2f}")

if __name__ == "__main__":
    main()
\end{lstlisting}
O método \texttt{groupby} junto do \texttt{sum()} nos permite somar todas as medições de um mesmo dia para criar uma agregação correta. Após isso, usamos \texttt{pd.merge} e \texttt{fillna(method='ffill')} para alinhar o ciclo com os dados de clima correspondentes e limpar as lacunas.
\end{solucao}

\vspace{2em}
\noindent\textbf{Créditos:} Este conteúdo foi originalmente criado pela \textit{The University of Helsinki MOOC Center} para o curso \textit{Data Analysis with Python}.
\end{document}
